\chapter{Ordinary Least Squares}
\label{cap:ols}
% ── index ───────────────────────────────────────────────────────
\index{ordinary least squares|textbf}
\index{OLS|see{ordinary least squares}}
\index{ols@\texttt{ols()}|textbf}
\index{omitted variable bias|textbf}
\index{BLUE|textbf}
\index{consistency}
\index{heteroscedasticity}
\index{robust standard errors}
\index{esttab@\texttt{esttab()}|textbf}

% ─────────────────────────────────────────────────────────────────────────────
\section{The model}

Consider the linear regression model:

\[
  y = X\beta + \varepsilon
\]

where $y$ is an $n \times 1$ vector of observations of the dependent variable,
$X$ is an $n \times k$ matrix of regressors (the first column is
identically 1, corresponding to the intercept), $\beta$ is a
$k \times 1$ vector of unknown parameters, and $\varepsilon$ is an
$n \times 1$ vector of errors.

\subsection*{Gauss-Markov Assumptions}

\begin{enumerate}
  \item \textbf{Linearity:} $y = X\beta + \varepsilon$.
  \item \textbf{Full rank:} $\mathrm{rank}(X) = k$ (no perfect multicollinearity).
  \item \textbf{Strict exogeneity:} $E[\varepsilon \mid X] = 0$.
  \item \textbf{Homoscedasticity:} $\mathrm{Var}[\varepsilon \mid X] = \sigma^2 I_n$.
  \item \textbf{No autocorrelation:} $\mathrm{Cov}[\varepsilon_i, \varepsilon_j \mid X] = 0$ for $i \neq j$.
\end{enumerate}

Under A1--A3 OLS is unbiased. Under A1--A5 it is BLUE (Gauss-Markov Theorem).

% ─────────────────────────────────────────────────────────────────────────────
\section{The estimator}

The OLS estimator minimizes the sum of squared residuals:

\[
  \hat{\beta} = \arg\min_{\beta}\, (y - X\beta)'(y - X\beta)
\]

The first-order condition $X'(y - X\beta) = 0$ --- the \emph{normal equations} ---
has an analytical solution:

\[
  \boxed{\hat{\beta} = (X'X)^{-1}X'y}
\]

The residuals are $\hat{e} = y - X\hat{\beta} = M_X y$, where
$M_X = I - X(X'X)^{-1}X'$ is the residual maker matrix (symmetric and idempotent).

% ─────────────────────────────────────────────────────────────────────────────
\section{Properties}

Under A1--A3:
\[
  E[\hat{\beta}] = \beta
\]

Under A1--A5, the conditional variance of the estimator is:
\[
  \mathrm{Var}[\hat{\beta} \mid X] = \sigma^2 (X'X)^{-1}
\]

The unbiased estimator of the error variance is:
\[
  s^2 = \frac{\hat{e}'\hat{e}}{n - k}
\]

and therefore $\widehat{\mathrm{Var}}[\hat{\beta}] = s^2 (X'X)^{-1}$.

% ─────────────────────────────────────────────────────────────────────────────
\section{DGP Verification}

The most direct proof that the estimator is correct is to define a DGP
(\emph{data-generating process}) with known parameters, generate data
from it and verify that OLS recovers them.

Choose the true parameters:
\[
  \beta_0 = 2{,}5 \qquad \beta_1 = 1{,}8 \qquad \beta_2 = -0{,}7
\]

and generate $n = 1000$ observations \emph{without} an error term ($\varepsilon = 0$):
\[
  y_i = 2{,}5 + 1{,}8\,x_{1i} - 0{,}7\,x_{2i}
\]

With $\varepsilon = 0$ the Gauss-Markov assumptions are satisfied exactly,
and OLS must recover $(\beta_0, \beta_1, \beta_2)$ to machine precision.

\begin{lstlisting}[caption={Noise-free DGP --- numerical verification}]
seed(42)
input df
id
1
end
for i in 1..=10 { df = append(df, df) }
df |> filter(_n <= 1000)
generate df x1 = rnormal()
generate df x2 = rnormal()
generate df y  = 2.5 + 1.8*x1 - 0.7*x2
ols(y ~ x1 + x2, df)
\end{lstlisting}

\noindent Actual interpreter output:

\begin{lstlisting}[language={}, basicstyle=\ttfamily\scriptsize,
  frame=none, numbers=none, backgroundcolor=\color{codbg},
  xleftmargin=0pt, xrightmargin=0pt, breaklines=false]
=========================== OLS Regression Results ===========================
Dep. Variable:                     y || R-squared:                    1.0000
Model:                           OLS || Adj. R-squared:               1.0000
Covariance Type:          Non-Robust || F-statistic:                  0.0000
No. Observations:               1000 || Prob (F-statistic):         1.0000e0
Df Residuals:                    997 || Log-Likelihood:           31513.9419
AIC:                     -63021.8837 || BIC:                     -63007.1605

------------------------------------------------------------------------------
Variable   |       coef |    std err |        t |    P>|t| | [0.025  0.975]
------------------------------------------------------------------------------
const      |     2.5000 |     0.0000 | 6.17e+15 |    0.000 |  2.5000  2.5000
x1         |     1.8000 |     0.0000 | 3.36e+15 |    0.000 |  1.8000  1.8000
x2         |    -0.7000 |     0.0000 |-1.28e+15 |    0.000 | -0.7000 -0.7000
==============================================================================
\end{lstlisting}

OLS recovers exactly $(\hat\beta_0, \hat\beta_1, \hat\beta_2) = (2{,}5;\; 1{,}8;\; -0{,}7)$.
The zero standard error and $R^2 = 1$ are expected: without noise, the data lie
exactly on the hyperplane defined by $\beta$, and the estimator is numerically
exact to floating-point precision.

\begin{nota}
  The $F$-statistic appears as \texttt{0.0} and the p-value as \texttt{1.0}
  in this case --- a numerical artifact of $0/0$ when $\mathrm{SSR} = 0$.
  The $t$-statistics are on the order of $10^{15}$ for the same reason: finite
  numerator divided by a machine-precision denominator.
  In the presence of noise ($\varepsilon \neq 0$) all statistics take
  finite and interpretable values.
\end{nota}

\subsection*{Misspecification effect}

The same DGP allows us to verify what happens when $x_2$ is omitted.
Two models are estimated and compared with \lstinline{esttab}:

\begin{lstlisting}[caption={Correct model vs.\ omitted $x_2$}]
let m1 = ols(y ~ x1 + x2, df)
let m2 = ols(y ~ x1, df)
esttab(m1, m2)
\end{lstlisting}

\begin{lstlisting}[language={}, basicstyle=\ttfamily\small,
  frame=none, numbers=none, backgroundcolor=\color{codbg},
  xleftmargin=0pt, xrightmargin=0pt, breaklines=false]
                            (1)              (2)
                            OLS              OLS
────────────────────────────────────────────────
x1                    1.8000***        1.7422***
                       (0.0000)         (0.0224)
x2                   -0.7000***
                       (0.0000)
Constant              2.5000***        2.1810***
                       (0.0000)         (0.0129)
────────────────────────────────────────────────
N                          1000             1000
R²                       1.0000           0.8585
Adj. R²                  1.0000           0.8583
────────────────────────────────────────────────
* p<0.10  ** p<0.05  *** p<0.01
\end{lstlisting}

\noindent The coefficient on $x_1$ barely changes (1.8000 $\to$ 1.7422) because
$x_1$ and $x_2$ are generated independently --- the sample correlation is
close to zero. The constant, however, is distorted from 2.5000 to 2.1810 because
it absorbs $\hat\beta_2\,\bar{x}_2$, and the sample mean of $x_2$ is not
exactly zero even with $n=1000$.

\begin{nota}
  Omitted variable bias (OVB) \emph{requires correlation} between the included
  and the omitted regressor. Without correlation, omission only inflates standard
  errors and reduces $R^2$ (loss of efficiency, not of consistency). For the
  formal treatment of endogeneity and instruments, see
  Chapter~\ref{cap:iv}.
\end{nota}

% ─────────────────────────────────────────────────────────────────────────────
\section{Standard errors}

\subsection*{Classical}

Assumes homoscedasticity. The estimated covariance matrix is
$s^2(X'X)^{-1}$ and the standard error of coefficient $j$ is
$\mathrm{se}_j = s\,\sqrt{[(X'X)^{-1}]_{jj}}$.

\subsection*{Heteroscedasticity-robust}

When $\mathrm{Var}[\varepsilon_i \mid x_i] = \sigma_i^2$ (heteroscedasticity),
the sandwich estimator provides valid inference.

\textbf{HC1} (White with degrees-of-freedom correction):
\[
  \hat{V}_{\mathrm{HC1}} = \frac{n}{n-k}\,(X'X)^{-1}
    \!\left(\sum_{i=1}^{n} \hat{e}_i^2\, x_i x_i'\right)\!(X'X)^{-1}
\]

\textbf{HC3} (more conservative, recommended in small samples):
\[
  \hat{V}_{\mathrm{HC3}} = (X'X)^{-1}
    \!\left(\sum_{i=1}^{n} \frac{\hat{e}_i^2}{(1 - h_{ii})^2}\, x_i x_i'\right)\!(X'X)^{-1}
\]

where $h_{ii} = x_i'(X'X)^{-1}x_i$ is the \emph{leverage} of observation $i$.

\subsection*{Clustered}

When errors are correlated within groups $g = 1,\ldots,G$:
\[
  \hat{V}_{\mathrm{cl}} = \frac{G(n-1)}{(G-1)(n-k)}\,
    (X'X)^{-1}\!\left(\sum_{g=1}^{G} X_g'\hat{e}_g\hat{e}_g'X_g\right)\!(X'X)^{-1}
\]

% ─────────────────────────────────────────────────────────────────────────────
\section{Inference}

\subsection*{$t$-test on a single coefficient}

Under H0: $\beta_j = \beta_j^0$:
\[
  t_j = \frac{\hat{\beta}_j - \beta_j^0}{\mathrm{se}(\hat{\beta}_j)}
  \;\xrightarrow{d}\; t_{n-k}
\]

\subsection*{$F$-test on linear restrictions}

For $R\beta = r$ with $q$ restrictions:
\[
  F = \frac{(R\hat{\beta} - r)'
       \bigl[R\,(X'X)^{-1}R'\bigr]^{-1}
       (R\hat{\beta} - r)}{q\,s^2}
  \;\xrightarrow{d}\; F_{q,\,n-k}
\]

\begin{nota}
  With robust or clustered errors, $\hat{V}$ replaces $s^2(X'X)^{-1}$
  in the computation of the $F$-test.
\end{nota}

% ─────────────────────────────────────────────────────────────────────────────
\section{Syntax}

\begin{lstlisting}
ols(formula, df)
ols(formula, df, cov=type)
ols(formula, df, cov=cluster, cluster=column)
ols(formula, df, if = condition)
\end{lstlisting}

Alias: \lstinline{reg()} is equivalent to \lstinline{ols()}.

\section{Formula}
\label{sec:ols-formula}
\index{formula!syntax}
\index{formula!transformations}
\index{formula!interactions}

The formula follows the standard notation \texttt{y \textasciitilde{} x1 + x2 + \ldots}:

\begin{lstlisting}[caption={Basic OLS}]
input df
Y X
10 2
12 3
8  1
15 5
11 2
14 4
end

ols(Y ~ X, df)
\end{lstlisting}

The intercept is included automatically. For multiple variables:

\begin{lstlisting}
ols(Y ~ X1 + X2 + X3, df)
\end{lstlisting}

The formula can be passed as a string:

\begin{lstlisting}
let formula = "Y ~ X1 + X2"
ols(formula, df)
\end{lstlisting}

\subsection{Transformations}
\label{subsec:formula-transformations}

Any function supported by the language can be used directly in the RHS
of the formula --- the interpreter evaluates the expression \emph{before}
passing the data to the estimator, so Greeners receives only flat numeric
columns.

\begin{lstlisting}[caption={Transformed regressors}]
// log of a variable
ols(Y ~ log(X), df)

// square root
ols(Y ~ sqrt(X), df)

// multiple transformations
ols(Y ~ log(K) + log(L), df)

// quadratic term (equivalent to I(X^2))
ols(Y ~ X + X^2, df)
\end{lstlisting}

The coefficient name in the output table reflects the original expression:
\texttt{log(K)}, \texttt{sqrt(X)}, etc.

\subsection{Interactions}
\label{subsec:formula-interactions}

The \texttt{:} operator creates an interaction term (element-wise product)
between two regressors. Each side can be an arbitrary expression:

\begin{lstlisting}[caption={Interactions}]
// interaction between simple variables
ols(Y ~ X1 + X2 + X1:X2, df)

// interaction of transformations --- log(K)*log(L)
ols(Y ~ log(K) + log(L) + log(K):log(L), df)

// interaction term only
ols(Y ~ log(K):log(L), df)
\end{lstlisting}

The \texttt{*} operator is syntactic sugar for full expansion
(\texttt{x1*x2} $\equiv$ \texttt{x1 + x2 + x1:x2}):

\begin{lstlisting}
// equivalent to: ols(Y ~ X1 + X2 + X1:X2, df)
ols(Y ~ X1 * X2, df)
\end{lstlisting}

\subsection{\texttt{I(expr)} --- arbitrary arithmetic expressions}
\label{subsec:formula-I}

The function \texttt{I(\ldots)} isolates a compound arithmetic expression
inside the formula, protecting the operators \texttt{+}, \texttt{-} and
\texttt{*} from the special meaning they carry in formula grammar:

\begin{lstlisting}[caption={Arbitrary terms with \texttt{I()}}]
// square of X
ols(Y ~ X + I(X^2), df)

// interaction term built explicitly
ols(Y ~ X1 + X2 + I(X1 * X2), df)

// linear combination of variables
ols(Y ~ I(X1 + X2), df)
\end{lstlisting}

\begin{nota}
  \texttt{I(X\textasciicircum{}2)} and \texttt{X\textasciicircum{}2} are equivalent when
  \texttt{X\textasciicircum{}2} appears as a standalone term (outside interactions and
  top-level additions). Use \texttt{I(\ldots)} to eliminate any ambiguity
  or when the expression contains \texttt{+} or \texttt{*}.
\end{nota}

\subsection{Categorical variable: \texttt{C()}}
\label{subsec:formula-categorical}

\texttt{C(var)} instructs the estimator to treat \texttt{var} as
categorical, automatically generating dummies (reference category = first
in lexicographic order):

\begin{lstlisting}
ols(Y ~ X + C(region), df)
\end{lstlisting}

\section{Standard error options}

\begin{center}
\begin{tabular}{ll}
\toprule
\textbf{Option} & \textbf{Formula} \\
\midrule
(default)                           & $s^2(X'X)^{-1}$ \\
\texttt{cov=robust}                 & HC1 (White sandwich) \\
\texttt{cov=HC1}                    & same \\
\texttt{cov=HC3}                    & HC3 (leverage-adjusted) \\
\texttt{cov=cluster, cluster=col}   & clustered by \texttt{col} \\
\bottomrule
\end{tabular}
\end{center}

\begin{lstlisting}[caption={Standard error variants}]
ols(Y ~ X, df)                           // classical
ols(Y ~ X, df, cov=robust)              // HC1
ols(Y ~ X, df, cov=HC3)                 // HC3
ols(Y ~ X, df, cov=cluster, cluster=id) // clustered
\end{lstlisting}

\section{Conditional sample}

The \lstinline{if =} option restricts estimation to a subset without modifying the DataFrame:

\begin{lstlisting}
ols(Y ~ X, df, if = group == 1)
\end{lstlisting}

\section{Capturing the result}

\begin{lstlisting}
let m = ols(Y ~ X, df)
display m
\end{lstlisting}

When assigned, \lstinline{ols()} does not print the table --- it only stores the model.

\section{Post-estimation}

\subsection{\texttt{predict}}

\begin{lstlisting}
let m = ols(Y ~ X, df)

// $\hat{y} = X\hat{\beta}$
let yhat = predict(m, "xb")
// $\hat{e} = y - X\hat{\beta}$
let ehat = predict(m, "residuals")
\end{lstlisting}

Accepted aliases: \texttt{"fitted"} for fitted values; \texttt{"resid"} or \texttt{"e"} for residuals.

\subsection{\texttt{vif}}

The VIF of regressor $j$ is $\mathrm{VIF}_j = 1/(1 - R_j^2)$, where $R_j^2$ is the
$R^2$ of the regression of $x_j$ on the remaining regressors. Values above 10
indicate severe multicollinearity.

\begin{lstlisting}
let m = ols(Y ~ X1 + X2 + X3, df)
vif(m)
\end{lstlisting}

\subsection{\texttt{test}}

$F$-test of linear restrictions $H_0: R\beta = r$:

\begin{lstlisting}
let m = ols(Y ~ X1 + X2 + X3, df)
test(m, X2 = 0, X3 = 0)    // H0: X2 = X3 = 0
\end{lstlisting}

\subsection{\texttt{lincom}}

Linear combination $c'\hat{\beta}$ with standard error $\sqrt{c'\hat{V}c}$:

\begin{lstlisting}
let m = ols(Y ~ X1 + X2, df)
lincom(m, X1 + X2)
\end{lstlisting}

\subsection{\texttt{margins}}

Average marginal effects. In linear OLS $\partial E[y]/\partial x_j = \beta_j$,
equivalent to the coefficients.

\begin{lstlisting}
let m = ols(Y ~ X, df)
margins(m)
\end{lstlisting}

\subsection{Specification diagnostics}

\subsubsection*{Ramsey RESET}

Tests specification error by adding powers of $\hat{y}$ as regressors:

\begin{lstlisting}
reset(m)              // default power 3
reset(m, power=2)     // only $\hat{y}^2$
\end{lstlisting}

\subsubsection*{CUSUM --- structural stability}

CUSUM test (Brown, Durbin, Evans 1975) using recursive residuals. Checks
whether the cumulative sum of recursive residuals stays within 5\%
significance bounds:

\begin{lstlisting}
cusumtest(m)
\end{lstlisting}

Rejects $H_0$ (stability) if the cumulative sum crosses the bounds ---
indicates a structural break in the sample period.

\subsection{Residual diagnostics}

\subsubsection*{Heteroskedasticity}

\begin{lstlisting}
test(m, "white")      // White (1980) --- general
test(m, "bp")         // Breusch-Pagan --- linear form
gqtest(m, split=0.2)  // Goldfeld-Quandt --- variance comparison
\end{lstlisting}

The Goldfeld-Quandt splits the sample into two groups (omitting the central
fraction \texttt{split}) and compares residual variances via an $F$ statistic.

\subsubsection*{Autocorrelation}

\begin{lstlisting}
test(m, "dw")         // Durbin-Watson --- AR(1)
bgodfrey(m, lags=4)   // Breusch-Godfrey --- AR(p) of arbitrary order
\end{lstlisting}

\subsubsection*{ACF and PACF of residuals}

\begin{lstlisting}
let a = acf(m, lags=10)    // ACF of residuals
let p = pacf(m, lags=10)   // PACF of residuals
ljungbox(m, lags=10)       // joint autocorrelation test
\end{lstlisting}

\texttt{acf} and \texttt{pacf} return lists of length \texttt{lags + 1}
(element 0 is always 1.0). They also accept \texttt{acf(df, var, lags=20)}
for raw series. For ASCII visualization, use \texttt{acfplot} and
\texttt{pacfplot}.

\subsubsection*{Normality}

\begin{lstlisting}
jb(m)                 // Jarque-Bera (skewness + kurtosis)
swilk(df, residuals)  // Shapiro-Wilk
\end{lstlisting}

\subsection{Influence diagnostics}

\begin{lstlisting}
influence(m)          // DFBETAS, DFFITS, Cook's D, leverage
leverage(m)           // hat matrix diagonal only
\end{lstlisting}

\subsection{Model comparison}

\begin{lstlisting}
estat(m1, m2, m3)              // AIC/BIC table with Akaike weights
let w = akaike_weights(m1, m2) // weights as dict {model: weight}
\end{lstlisting}

\texttt{akaike\_weights} returns a dictionary with Akaike weights
$w_i = \frac{e^{-\Delta_i/2}}{\sum_j e^{-\Delta_j/2}}$, where $\Delta_i =
\mathrm{AIC}_i - \min_j \mathrm{AIC}_j$. Supports OLS, logit/probit, Poisson,
NegBin, Tobit, Ordered, Mixed, and Zero-Inflated.

\subsubsection*{Likelihood-ratio test}

Compares two nested models --- the restricted is a special case of the
unrestricted:

\begin{lstlisting}
let m1 = ols(Y ~ X1, df)
let m2 = ols(Y ~ X1 + X2, df)
lrtest(m1, m2)    // H0: m1 is adequate (X2 = 0)
\end{lstlisting}

The statistic is $\mathrm{LR} = -2(\ln L_r - \ln L_u) \sim \chi^2(\mathrm{df})$,
where $\mathrm{df} = k_u - k_r$. Rejecting $H_0$ indicates the unrestricted
model fits significantly better. Works with any estimator that reports a
log-likelihood: OLS, logit/probit, Poisson, NegBin, Tobit, Ordered, Mixed,
ZI, GLM, GARCH, ARIMA.

\section{Results table with \texttt{esttab}}

\begin{lstlisting}[caption={Comparing specifications}]
let m1 = ols(Y ~ X, df)
let m2 = ols(Y ~ X, df, cov=robust)
let m3 = ols(Y ~ X1 + X2, df)

esttab(m1, m2, m3)
\end{lstlisting}

\begin{lstlisting}
eststo(ols(Y ~ X, df))
eststo(ols(Y ~ X1 + X2, df))
esttab()
\end{lstlisting}

\section{Bootstrap}

Bootstrap standard errors --- a nonparametric alternative:

\begin{lstlisting}
bootstrap(ols, Y ~ X, df, n=1000)
bootse(m, n=500)
\end{lstlisting}
